\documentclass[11pt,a4paper]{article}
% main (standalone preamble for editing)
% vim: nowrap: spell:
% ══════════════════════════════════════════════════════════════════════════════
% Speed Maths --- shared preamble
% Included by every sheet and answer file via:  % main (standalone preamble for editing)
% vim: nowrap: spell:
% ══════════════════════════════════════════════════════════════════════════════
% Speed Maths --- shared preamble
% Included by every sheet and answer file via:  % main (standalone preamble for editing)
% vim: nowrap: spell:
% ══════════════════════════════════════════════════════════════════════════════
% Speed Maths --- shared preamble
% Included by every sheet and answer file via:  \input{../../shared/preamble}
% Editing THIS file changes the look of every sheet/answer across every pillar.
% ══════════════════════════════════════════════════════════════════════════════

\usepackage[a4paper, top=25mm, bottom=25mm, left=22mm, right=22mm]{geometry}
\usepackage{amsmath, amssymb}
\usepackage{enumitem}
\usepackage{booktabs}
\usepackage{array}
\usepackage{parskip}
\usepackage{microtype}
\usepackage{fancyhdr}
\usepackage{titlesec}
\usepackage{mdframed}
\usepackage{xcolor}
\usepackage[hidelinks]{hyperref}
\usepackage{graphicx}
\usepackage{tikz}
\usepackage{pgfplots}
\pgfplotsset{compat=1.18}

% ── Colours ───────────────────────────────────────────────────────────────────
\definecolor{darkgray}{gray}{0.15}
\definecolor{medgray}{gray}{0.45}
\definecolor{lightgray}{gray}{0.93}
\definecolor{boxgray}{gray}{0.88}

% ── Section formatting ──────────────────────────────────────────────────────────
\titleformat{\section}[block]
  {\normalfont\large\bfseries}{}
  {0pt}{}[\vspace{2pt}\hrule\vspace{6pt}]
\titlespacing*{\section}{0pt}{14pt}{4pt}

% ── List formatting ─────────────────────────────────────────────────────────────
\setlist[enumerate,1]{label=\textbf{\arabic*.}, leftmargin=2em, itemsep=10pt}

% ── Branding constants (edit here, not per-file) ────────────────────────────────
\newcommand{\SpeedExamLine}{\textbf{Speed Maths:} Competition \& Exam Prep}
\newcommand{\SpeedCredit}{Series by \href{https://www.linkedin.com/in/cerealdev/}{David Akinyele-Aje}}

% ── Header / footer ──────────────────────────────────────────────────────────────
% Usage (once, after \input{../../shared/preamble} and before \begin{document}):
%   \SpeedHeader{Algebra}{3}
% First arg  = pillar name shown as "Speed <Pillar> --- Daily Drill"
% Second arg = worksheet number
\pagestyle{fancy}
\newcommand{\SpeedHeader}[2]{%
  \fancyhf{}%
  \renewcommand{\headrulewidth}{0.4pt}%
  \setlength{\headheight}{14pt}%
  \fancyhead[L]{\small\textbf{Speed #1 --- Daily Drill}}%
  \fancyhead[R]{\small Worksheet \##2}%
  \fancyfoot[L]{\small \SpeedExamLine}%
  \fancyfoot[C]{\small\thepage}%
  \fancyfoot[R]{\small \SpeedCredit}%
  \renewcommand{\footrulewidth}{0.4pt}%
}

% ── Title block (used at the top of every sheet AND answer file) ───────────────
% Usage: \SpeedTitleBlock{Daily Algebra Drill \#3}{Jane Doe}
% %% AGENTS: When generating a new sheet, ask the human contributor for the name they want credited in argument 2.
% For answer files, pass the "--- Answers \& Investigations" suffix in the arg.
\newcommand{\SpeedTitleBlock}[2]{%
  \begin{center}
    {\LARGE\bfseries #1}\\[4pt]
    {\normalsize\itshape Sheet by #2}\\[6pt]
    \hrule\vspace{4pt}
    \begin{tabular}{llcc}
      \toprule
      \textbf{Section} & \textbf{Topic} & \textbf{Questions} & \textbf{Target time}\\
      \midrule
      A & Rapid Recognition          & 10 & 2 min 30 sec \\
      B & Manipulation Drills        & 10 & 8 min        \\
      C & Substitution \& Structure  & 8  & 10 min       \\
      D & Challenge                  & 5  & 15 min       \\
      \bottomrule
    \end{tabular}
    \vspace{4pt}
    \hrule
  \end{center}
}

% ── Sheet metadata: topic + the day's new tools ────────────────────────────────
% Usage (immediately after \SpeedTitleBlock, sheets only):
%   \SpeedMeta{Expanding, factorising, and the identity toolkit}
%             {difference of squares, sum/product form, completing the square}
%
% Arg 1 = the sheet's topic in one line. The website lists this instead of
%         "Sheet 03", so write the thing a reader would actually search for.
% Arg 2 = comma-separated, at most five: the tools this sheet introduces that
%         earlier sheets in the pillar did not. These become the website's tags.
%         Leave out anything the reader already met — the tags are meant to say
%         what is new today, not to summarise the sheet.
%
% Both are parsed straight out of this file by tools/build_website.py, so the
% sheet stays the single source of truth and the site cannot drift from it.
%
% This renders NOTHING. It is a declaration for the build script, not a piece of
% the sheet's layout: adding it to a sheet must not change that sheet's PDF, so
% the 25 existing sheets can be annotated without a single one of them being
% reissued. If the toolkit should also appear on the page, that is a separate
% decision about the sheet design — make it deliberately, here, not by leaving
% this macro printing something nobody asked it to.
\newcommand{\SpeedMeta}[2]{}

% ── Answer / Method / Investigate-further boxes (answer files only) ────────────
\newmdenv[
  backgroundcolor=lightgray,
  linecolor=medgray,
  linewidth=0.5pt,
  leftmargin=0pt, rightmargin=0pt,
  innerleftmargin=8pt, innerrightmargin=8pt,
  innertopmargin=5pt, innerbottommargin=5pt,
  skipabove=4pt, skipbelow=4pt
]{ansbox}

\newmdenv[
  backgroundcolor=white,
  linecolor=darkgray,
  linewidth=0.8pt,
  leftline=true, rightline=false, topline=false, bottomline=false,
  leftmargin=0pt, rightmargin=0pt,
  innerleftmargin=8pt, innerrightmargin=6pt,
  innertopmargin=4pt, innerbottommargin=4pt,
  skipabove=4pt, skipbelow=6pt
]{invbox}

\newcommand{\ans}[1]{%
  \begin{ansbox}
    \textbf{Answer:} #1
  \end{ansbox}}

\newcommand{\method}[1]{%
  {\small\textbf{Method:} #1}\par}

\newcommand{\inv}[1]{%
  \begin{invbox}
    \textbf{\small Investigate further:} {\small #1}
  \end{invbox}}

% ── Misc helpers ────────────────────────────────────────────────────────────────
\newcommand{\qn}[1]{\textbf{#1.}\;}

% Mistake-linking: attach a Top 5 Patterns / Common Traps bullet to the question
% label(s) in this sheet where it actually shows up, e.g. \seealso{B6, D2}
\newcommand{\seealso}[1]{\hfill\textit{\footnotesize(see #1)}}

% ── Graph options macro ─────────────────────────────────────────────────────────
% Usage: \GraphOptions{tikz code for A}{tikz code for B}{tikz code for C}{tikz code for D}
\newcommand{\GraphOptions}[4]{%
  \begin{center}
    \begin{tabular}{cc}
      \textbf{A)} & \textbf{B)} \\
      #1 & #2 \\[10pt]
      \textbf{C)} & \textbf{D)} \\
      #3 & #4
    \end{tabular}
  \end{center}
}

% ── Closing motivational rule + cross-promo footer (sheets only) ────────────────
% Usage: \SpeedClosing{"Quote text."}
\newcommand{\SpeedClosing}[1]{%
  \vspace{20pt}
  \begin{center}
  \rule{0.6\textwidth}{0.4pt}\\[4pt]
  {\small\itshape #1}\\[4pt]
  \rule{0.6\textwidth}{0.4pt}
  \end{center}
  \begin{center}
  {\small Found a mistake or want to contribute a sheet?}\\
  {\small Visit the open-source project at \href{https://github.com/The-CerealDev/speed-maths}{github.com/The-CerealDev/speed-maths}}
  \end{center}
}

% Editing THIS file changes the look of every sheet/answer across every pillar.
% ══════════════════════════════════════════════════════════════════════════════

\usepackage[a4paper, top=25mm, bottom=25mm, left=22mm, right=22mm]{geometry}
\usepackage{amsmath, amssymb}
\usepackage{enumitem}
\usepackage{booktabs}
\usepackage{array}
\usepackage{parskip}
\usepackage{microtype}
\usepackage{fancyhdr}
\usepackage{titlesec}
\usepackage{mdframed}
\usepackage{xcolor}
\usepackage[hidelinks]{hyperref}
\usepackage{graphicx}
\usepackage{tikz}
\usepackage{pgfplots}
\pgfplotsset{compat=1.18}

% ── Colours ───────────────────────────────────────────────────────────────────
\definecolor{darkgray}{gray}{0.15}
\definecolor{medgray}{gray}{0.45}
\definecolor{lightgray}{gray}{0.93}
\definecolor{boxgray}{gray}{0.88}

% ── Section formatting ──────────────────────────────────────────────────────────
\titleformat{\section}[block]
  {\normalfont\large\bfseries}{}
  {0pt}{}[\vspace{2pt}\hrule\vspace{6pt}]
\titlespacing*{\section}{0pt}{14pt}{4pt}

% ── List formatting ─────────────────────────────────────────────────────────────
\setlist[enumerate,1]{label=\textbf{\arabic*.}, leftmargin=2em, itemsep=10pt}

% ── Branding constants (edit here, not per-file) ────────────────────────────────
\newcommand{\SpeedExamLine}{\textbf{Speed Maths:} Competition \& Exam Prep}
\newcommand{\SpeedCredit}{Series by \href{https://www.linkedin.com/in/cerealdev/}{David Akinyele-Aje}}

% ── Header / footer ──────────────────────────────────────────────────────────────
% Usage (once, after % main (standalone preamble for editing)
% vim: nowrap: spell:
% ══════════════════════════════════════════════════════════════════════════════
% Speed Maths --- shared preamble
% Included by every sheet and answer file via:  \input{../../shared/preamble}
% Editing THIS file changes the look of every sheet/answer across every pillar.
% ══════════════════════════════════════════════════════════════════════════════

\usepackage[a4paper, top=25mm, bottom=25mm, left=22mm, right=22mm]{geometry}
\usepackage{amsmath, amssymb}
\usepackage{enumitem}
\usepackage{booktabs}
\usepackage{array}
\usepackage{parskip}
\usepackage{microtype}
\usepackage{fancyhdr}
\usepackage{titlesec}
\usepackage{mdframed}
\usepackage{xcolor}
\usepackage[hidelinks]{hyperref}
\usepackage{graphicx}
\usepackage{tikz}
\usepackage{pgfplots}
\pgfplotsset{compat=1.18}

% ── Colours ───────────────────────────────────────────────────────────────────
\definecolor{darkgray}{gray}{0.15}
\definecolor{medgray}{gray}{0.45}
\definecolor{lightgray}{gray}{0.93}
\definecolor{boxgray}{gray}{0.88}

% ── Section formatting ──────────────────────────────────────────────────────────
\titleformat{\section}[block]
  {\normalfont\large\bfseries}{}
  {0pt}{}[\vspace{2pt}\hrule\vspace{6pt}]
\titlespacing*{\section}{0pt}{14pt}{4pt}

% ── List formatting ─────────────────────────────────────────────────────────────
\setlist[enumerate,1]{label=\textbf{\arabic*.}, leftmargin=2em, itemsep=10pt}

% ── Branding constants (edit here, not per-file) ────────────────────────────────
\newcommand{\SpeedExamLine}{\textbf{Speed Maths:} Competition \& Exam Prep}
\newcommand{\SpeedCredit}{Series by \href{https://www.linkedin.com/in/cerealdev/}{David Akinyele-Aje}}

% ── Header / footer ──────────────────────────────────────────────────────────────
% Usage (once, after \input{../../shared/preamble} and before \begin{document}):
%   \SpeedHeader{Algebra}{3}
% First arg  = pillar name shown as "Speed <Pillar> --- Daily Drill"
% Second arg = worksheet number
\pagestyle{fancy}
\newcommand{\SpeedHeader}[2]{%
  \fancyhf{}%
  \renewcommand{\headrulewidth}{0.4pt}%
  \setlength{\headheight}{14pt}%
  \fancyhead[L]{\small\textbf{Speed #1 --- Daily Drill}}%
  \fancyhead[R]{\small Worksheet \##2}%
  \fancyfoot[L]{\small \SpeedExamLine}%
  \fancyfoot[C]{\small\thepage}%
  \fancyfoot[R]{\small \SpeedCredit}%
  \renewcommand{\footrulewidth}{0.4pt}%
}

% ── Title block (used at the top of every sheet AND answer file) ───────────────
% Usage: \SpeedTitleBlock{Daily Algebra Drill \#3}{Jane Doe}
% %% AGENTS: When generating a new sheet, ask the human contributor for the name they want credited in argument 2.
% For answer files, pass the "--- Answers \& Investigations" suffix in the arg.
\newcommand{\SpeedTitleBlock}[2]{%
  \begin{center}
    {\LARGE\bfseries #1}\\[4pt]
    {\normalsize\itshape Sheet by #2}\\[6pt]
    \hrule\vspace{4pt}
    \begin{tabular}{llcc}
      \toprule
      \textbf{Section} & \textbf{Topic} & \textbf{Questions} & \textbf{Target time}\\
      \midrule
      A & Rapid Recognition          & 10 & 2 min 30 sec \\
      B & Manipulation Drills        & 10 & 8 min        \\
      C & Substitution \& Structure  & 8  & 10 min       \\
      D & Challenge                  & 5  & 15 min       \\
      \bottomrule
    \end{tabular}
    \vspace{4pt}
    \hrule
  \end{center}
}

% ── Sheet metadata: topic + the day's new tools ────────────────────────────────
% Usage (immediately after \SpeedTitleBlock, sheets only):
%   \SpeedMeta{Expanding, factorising, and the identity toolkit}
%             {difference of squares, sum/product form, completing the square}
%
% Arg 1 = the sheet's topic in one line. The website lists this instead of
%         "Sheet 03", so write the thing a reader would actually search for.
% Arg 2 = comma-separated, at most five: the tools this sheet introduces that
%         earlier sheets in the pillar did not. These become the website's tags.
%         Leave out anything the reader already met — the tags are meant to say
%         what is new today, not to summarise the sheet.
%
% Both are parsed straight out of this file by tools/build_website.py, so the
% sheet stays the single source of truth and the site cannot drift from it.
%
% This renders NOTHING. It is a declaration for the build script, not a piece of
% the sheet's layout: adding it to a sheet must not change that sheet's PDF, so
% the 25 existing sheets can be annotated without a single one of them being
% reissued. If the toolkit should also appear on the page, that is a separate
% decision about the sheet design — make it deliberately, here, not by leaving
% this macro printing something nobody asked it to.
\newcommand{\SpeedMeta}[2]{}

% ── Answer / Method / Investigate-further boxes (answer files only) ────────────
\newmdenv[
  backgroundcolor=lightgray,
  linecolor=medgray,
  linewidth=0.5pt,
  leftmargin=0pt, rightmargin=0pt,
  innerleftmargin=8pt, innerrightmargin=8pt,
  innertopmargin=5pt, innerbottommargin=5pt,
  skipabove=4pt, skipbelow=4pt
]{ansbox}

\newmdenv[
  backgroundcolor=white,
  linecolor=darkgray,
  linewidth=0.8pt,
  leftline=true, rightline=false, topline=false, bottomline=false,
  leftmargin=0pt, rightmargin=0pt,
  innerleftmargin=8pt, innerrightmargin=6pt,
  innertopmargin=4pt, innerbottommargin=4pt,
  skipabove=4pt, skipbelow=6pt
]{invbox}

\newcommand{\ans}[1]{%
  \begin{ansbox}
    \textbf{Answer:} #1
  \end{ansbox}}

\newcommand{\method}[1]{%
  {\small\textbf{Method:} #1}\par}

\newcommand{\inv}[1]{%
  \begin{invbox}
    \textbf{\small Investigate further:} {\small #1}
  \end{invbox}}

% ── Misc helpers ────────────────────────────────────────────────────────────────
\newcommand{\qn}[1]{\textbf{#1.}\;}

% Mistake-linking: attach a Top 5 Patterns / Common Traps bullet to the question
% label(s) in this sheet where it actually shows up, e.g. \seealso{B6, D2}
\newcommand{\seealso}[1]{\hfill\textit{\footnotesize(see #1)}}

% ── Graph options macro ─────────────────────────────────────────────────────────
% Usage: \GraphOptions{tikz code for A}{tikz code for B}{tikz code for C}{tikz code for D}
\newcommand{\GraphOptions}[4]{%
  \begin{center}
    \begin{tabular}{cc}
      \textbf{A)} & \textbf{B)} \\
      #1 & #2 \\[10pt]
      \textbf{C)} & \textbf{D)} \\
      #3 & #4
    \end{tabular}
  \end{center}
}

% ── Closing motivational rule + cross-promo footer (sheets only) ────────────────
% Usage: \SpeedClosing{"Quote text."}
\newcommand{\SpeedClosing}[1]{%
  \vspace{20pt}
  \begin{center}
  \rule{0.6\textwidth}{0.4pt}\\[4pt]
  {\small\itshape #1}\\[4pt]
  \rule{0.6\textwidth}{0.4pt}
  \end{center}
  \begin{center}
  {\small Found a mistake or want to contribute a sheet?}\\
  {\small Visit the open-source project at \href{https://github.com/The-CerealDev/speed-maths}{github.com/The-CerealDev/speed-maths}}
  \end{center}
}
 and before \begin{document}):
%   \SpeedHeader{Algebra}{3}
% First arg  = pillar name shown as "Speed <Pillar> --- Daily Drill"
% Second arg = worksheet number
\pagestyle{fancy}
\newcommand{\SpeedHeader}[2]{%
  \fancyhf{}%
  \renewcommand{\headrulewidth}{0.4pt}%
  \setlength{\headheight}{14pt}%
  \fancyhead[L]{\small\textbf{Speed #1 --- Daily Drill}}%
  \fancyhead[R]{\small Worksheet \##2}%
  \fancyfoot[L]{\small \SpeedExamLine}%
  \fancyfoot[C]{\small\thepage}%
  \fancyfoot[R]{\small \SpeedCredit}%
  \renewcommand{\footrulewidth}{0.4pt}%
}

% ── Title block (used at the top of every sheet AND answer file) ───────────────
% Usage: \SpeedTitleBlock{Daily Algebra Drill \#3}{Jane Doe}
% %% AGENTS: When generating a new sheet, ask the human contributor for the name they want credited in argument 2.
% For answer files, pass the "--- Answers \& Investigations" suffix in the arg.
\newcommand{\SpeedTitleBlock}[2]{%
  \begin{center}
    {\LARGE\bfseries #1}\\[4pt]
    {\normalsize\itshape Sheet by #2}\\[6pt]
    \hrule\vspace{4pt}
    \begin{tabular}{llcc}
      \toprule
      \textbf{Section} & \textbf{Topic} & \textbf{Questions} & \textbf{Target time}\\
      \midrule
      A & Rapid Recognition          & 10 & 2 min 30 sec \\
      B & Manipulation Drills        & 10 & 8 min        \\
      C & Substitution \& Structure  & 8  & 10 min       \\
      D & Challenge                  & 5  & 15 min       \\
      \bottomrule
    \end{tabular}
    \vspace{4pt}
    \hrule
  \end{center}
}

% ── Sheet metadata: topic + the day's new tools ────────────────────────────────
% Usage (immediately after \SpeedTitleBlock, sheets only):
%   \SpeedMeta{Expanding, factorising, and the identity toolkit}
%             {difference of squares, sum/product form, completing the square}
%
% Arg 1 = the sheet's topic in one line. The website lists this instead of
%         "Sheet 03", so write the thing a reader would actually search for.
% Arg 2 = comma-separated, at most five: the tools this sheet introduces that
%         earlier sheets in the pillar did not. These become the website's tags.
%         Leave out anything the reader already met — the tags are meant to say
%         what is new today, not to summarise the sheet.
%
% Both are parsed straight out of this file by tools/build_website.py, so the
% sheet stays the single source of truth and the site cannot drift from it.
%
% This renders NOTHING. It is a declaration for the build script, not a piece of
% the sheet's layout: adding it to a sheet must not change that sheet's PDF, so
% the 25 existing sheets can be annotated without a single one of them being
% reissued. If the toolkit should also appear on the page, that is a separate
% decision about the sheet design — make it deliberately, here, not by leaving
% this macro printing something nobody asked it to.
\newcommand{\SpeedMeta}[2]{}

% ── Answer / Method / Investigate-further boxes (answer files only) ────────────
\newmdenv[
  backgroundcolor=lightgray,
  linecolor=medgray,
  linewidth=0.5pt,
  leftmargin=0pt, rightmargin=0pt,
  innerleftmargin=8pt, innerrightmargin=8pt,
  innertopmargin=5pt, innerbottommargin=5pt,
  skipabove=4pt, skipbelow=4pt
]{ansbox}

\newmdenv[
  backgroundcolor=white,
  linecolor=darkgray,
  linewidth=0.8pt,
  leftline=true, rightline=false, topline=false, bottomline=false,
  leftmargin=0pt, rightmargin=0pt,
  innerleftmargin=8pt, innerrightmargin=6pt,
  innertopmargin=4pt, innerbottommargin=4pt,
  skipabove=4pt, skipbelow=6pt
]{invbox}

\newcommand{\ans}[1]{%
  \begin{ansbox}
    \textbf{Answer:} #1
  \end{ansbox}}

\newcommand{\method}[1]{%
  {\small\textbf{Method:} #1}\par}

\newcommand{\inv}[1]{%
  \begin{invbox}
    \textbf{\small Investigate further:} {\small #1}
  \end{invbox}}

% ── Misc helpers ────────────────────────────────────────────────────────────────
\newcommand{\qn}[1]{\textbf{#1.}\;}

% Mistake-linking: attach a Top 5 Patterns / Common Traps bullet to the question
% label(s) in this sheet where it actually shows up, e.g. \seealso{B6, D2}
\newcommand{\seealso}[1]{\hfill\textit{\footnotesize(see #1)}}

% ── Graph options macro ─────────────────────────────────────────────────────────
% Usage: \GraphOptions{tikz code for A}{tikz code for B}{tikz code for C}{tikz code for D}
\newcommand{\GraphOptions}[4]{%
  \begin{center}
    \begin{tabular}{cc}
      \textbf{A)} & \textbf{B)} \\
      #1 & #2 \\[10pt]
      \textbf{C)} & \textbf{D)} \\
      #3 & #4
    \end{tabular}
  \end{center}
}

% ── Closing motivational rule + cross-promo footer (sheets only) ────────────────
% Usage: \SpeedClosing{"Quote text."}
\newcommand{\SpeedClosing}[1]{%
  \vspace{20pt}
  \begin{center}
  \rule{0.6\textwidth}{0.4pt}\\[4pt]
  {\small\itshape #1}\\[4pt]
  \rule{0.6\textwidth}{0.4pt}
  \end{center}
  \begin{center}
  {\small Found a mistake or want to contribute a sheet?}\\
  {\small Visit the open-source project at \href{https://github.com/The-CerealDev/speed-maths}{github.com/The-CerealDev/speed-maths}}
  \end{center}
}

% Editing THIS file changes the look of every sheet/answer across every pillar.
% ══════════════════════════════════════════════════════════════════════════════

\usepackage[a4paper, top=25mm, bottom=25mm, left=22mm, right=22mm]{geometry}
\usepackage{amsmath, amssymb}
\usepackage{enumitem}
\usepackage{booktabs}
\usepackage{array}
\usepackage{parskip}
\usepackage{microtype}
\usepackage{fancyhdr}
\usepackage{titlesec}
\usepackage{mdframed}
\usepackage{xcolor}
\usepackage[hidelinks]{hyperref}
\usepackage{graphicx}
\usepackage{tikz}
\usepackage{pgfplots}
\pgfplotsset{compat=1.18}

% ── Colours ───────────────────────────────────────────────────────────────────
\definecolor{darkgray}{gray}{0.15}
\definecolor{medgray}{gray}{0.45}
\definecolor{lightgray}{gray}{0.93}
\definecolor{boxgray}{gray}{0.88}

% ── Section formatting ──────────────────────────────────────────────────────────
\titleformat{\section}[block]
  {\normalfont\large\bfseries}{}
  {0pt}{}[\vspace{2pt}\hrule\vspace{6pt}]
\titlespacing*{\section}{0pt}{14pt}{4pt}

% ── List formatting ─────────────────────────────────────────────────────────────
\setlist[enumerate,1]{label=\textbf{\arabic*.}, leftmargin=2em, itemsep=10pt}

% ── Branding constants (edit here, not per-file) ────────────────────────────────
\newcommand{\SpeedExamLine}{\textbf{Speed Maths:} Competition \& Exam Prep}
\newcommand{\SpeedCredit}{Series by \href{https://www.linkedin.com/in/cerealdev/}{David Akinyele-Aje}}

% ── Header / footer ──────────────────────────────────────────────────────────────
% Usage (once, after % main (standalone preamble for editing)
% vim: nowrap: spell:
% ══════════════════════════════════════════════════════════════════════════════
% Speed Maths --- shared preamble
% Included by every sheet and answer file via:  % main (standalone preamble for editing)
% vim: nowrap: spell:
% ══════════════════════════════════════════════════════════════════════════════
% Speed Maths --- shared preamble
% Included by every sheet and answer file via:  \input{../../shared/preamble}
% Editing THIS file changes the look of every sheet/answer across every pillar.
% ══════════════════════════════════════════════════════════════════════════════

\usepackage[a4paper, top=25mm, bottom=25mm, left=22mm, right=22mm]{geometry}
\usepackage{amsmath, amssymb}
\usepackage{enumitem}
\usepackage{booktabs}
\usepackage{array}
\usepackage{parskip}
\usepackage{microtype}
\usepackage{fancyhdr}
\usepackage{titlesec}
\usepackage{mdframed}
\usepackage{xcolor}
\usepackage[hidelinks]{hyperref}
\usepackage{graphicx}
\usepackage{tikz}
\usepackage{pgfplots}
\pgfplotsset{compat=1.18}

% ── Colours ───────────────────────────────────────────────────────────────────
\definecolor{darkgray}{gray}{0.15}
\definecolor{medgray}{gray}{0.45}
\definecolor{lightgray}{gray}{0.93}
\definecolor{boxgray}{gray}{0.88}

% ── Section formatting ──────────────────────────────────────────────────────────
\titleformat{\section}[block]
  {\normalfont\large\bfseries}{}
  {0pt}{}[\vspace{2pt}\hrule\vspace{6pt}]
\titlespacing*{\section}{0pt}{14pt}{4pt}

% ── List formatting ─────────────────────────────────────────────────────────────
\setlist[enumerate,1]{label=\textbf{\arabic*.}, leftmargin=2em, itemsep=10pt}

% ── Branding constants (edit here, not per-file) ────────────────────────────────
\newcommand{\SpeedExamLine}{\textbf{Speed Maths:} Competition \& Exam Prep}
\newcommand{\SpeedCredit}{Series by \href{https://www.linkedin.com/in/cerealdev/}{David Akinyele-Aje}}

% ── Header / footer ──────────────────────────────────────────────────────────────
% Usage (once, after \input{../../shared/preamble} and before \begin{document}):
%   \SpeedHeader{Algebra}{3}
% First arg  = pillar name shown as "Speed <Pillar> --- Daily Drill"
% Second arg = worksheet number
\pagestyle{fancy}
\newcommand{\SpeedHeader}[2]{%
  \fancyhf{}%
  \renewcommand{\headrulewidth}{0.4pt}%
  \setlength{\headheight}{14pt}%
  \fancyhead[L]{\small\textbf{Speed #1 --- Daily Drill}}%
  \fancyhead[R]{\small Worksheet \##2}%
  \fancyfoot[L]{\small \SpeedExamLine}%
  \fancyfoot[C]{\small\thepage}%
  \fancyfoot[R]{\small \SpeedCredit}%
  \renewcommand{\footrulewidth}{0.4pt}%
}

% ── Title block (used at the top of every sheet AND answer file) ───────────────
% Usage: \SpeedTitleBlock{Daily Algebra Drill \#3}{Jane Doe}
% %% AGENTS: When generating a new sheet, ask the human contributor for the name they want credited in argument 2.
% For answer files, pass the "--- Answers \& Investigations" suffix in the arg.
\newcommand{\SpeedTitleBlock}[2]{%
  \begin{center}
    {\LARGE\bfseries #1}\\[4pt]
    {\normalsize\itshape Sheet by #2}\\[6pt]
    \hrule\vspace{4pt}
    \begin{tabular}{llcc}
      \toprule
      \textbf{Section} & \textbf{Topic} & \textbf{Questions} & \textbf{Target time}\\
      \midrule
      A & Rapid Recognition          & 10 & 2 min 30 sec \\
      B & Manipulation Drills        & 10 & 8 min        \\
      C & Substitution \& Structure  & 8  & 10 min       \\
      D & Challenge                  & 5  & 15 min       \\
      \bottomrule
    \end{tabular}
    \vspace{4pt}
    \hrule
  \end{center}
}

% ── Sheet metadata: topic + the day's new tools ────────────────────────────────
% Usage (immediately after \SpeedTitleBlock, sheets only):
%   \SpeedMeta{Expanding, factorising, and the identity toolkit}
%             {difference of squares, sum/product form, completing the square}
%
% Arg 1 = the sheet's topic in one line. The website lists this instead of
%         "Sheet 03", so write the thing a reader would actually search for.
% Arg 2 = comma-separated, at most five: the tools this sheet introduces that
%         earlier sheets in the pillar did not. These become the website's tags.
%         Leave out anything the reader already met — the tags are meant to say
%         what is new today, not to summarise the sheet.
%
% Both are parsed straight out of this file by tools/build_website.py, so the
% sheet stays the single source of truth and the site cannot drift from it.
%
% This renders NOTHING. It is a declaration for the build script, not a piece of
% the sheet's layout: adding it to a sheet must not change that sheet's PDF, so
% the 25 existing sheets can be annotated without a single one of them being
% reissued. If the toolkit should also appear on the page, that is a separate
% decision about the sheet design — make it deliberately, here, not by leaving
% this macro printing something nobody asked it to.
\newcommand{\SpeedMeta}[2]{}

% ── Answer / Method / Investigate-further boxes (answer files only) ────────────
\newmdenv[
  backgroundcolor=lightgray,
  linecolor=medgray,
  linewidth=0.5pt,
  leftmargin=0pt, rightmargin=0pt,
  innerleftmargin=8pt, innerrightmargin=8pt,
  innertopmargin=5pt, innerbottommargin=5pt,
  skipabove=4pt, skipbelow=4pt
]{ansbox}

\newmdenv[
  backgroundcolor=white,
  linecolor=darkgray,
  linewidth=0.8pt,
  leftline=true, rightline=false, topline=false, bottomline=false,
  leftmargin=0pt, rightmargin=0pt,
  innerleftmargin=8pt, innerrightmargin=6pt,
  innertopmargin=4pt, innerbottommargin=4pt,
  skipabove=4pt, skipbelow=6pt
]{invbox}

\newcommand{\ans}[1]{%
  \begin{ansbox}
    \textbf{Answer:} #1
  \end{ansbox}}

\newcommand{\method}[1]{%
  {\small\textbf{Method:} #1}\par}

\newcommand{\inv}[1]{%
  \begin{invbox}
    \textbf{\small Investigate further:} {\small #1}
  \end{invbox}}

% ── Misc helpers ────────────────────────────────────────────────────────────────
\newcommand{\qn}[1]{\textbf{#1.}\;}

% Mistake-linking: attach a Top 5 Patterns / Common Traps bullet to the question
% label(s) in this sheet where it actually shows up, e.g. \seealso{B6, D2}
\newcommand{\seealso}[1]{\hfill\textit{\footnotesize(see #1)}}

% ── Graph options macro ─────────────────────────────────────────────────────────
% Usage: \GraphOptions{tikz code for A}{tikz code for B}{tikz code for C}{tikz code for D}
\newcommand{\GraphOptions}[4]{%
  \begin{center}
    \begin{tabular}{cc}
      \textbf{A)} & \textbf{B)} \\
      #1 & #2 \\[10pt]
      \textbf{C)} & \textbf{D)} \\
      #3 & #4
    \end{tabular}
  \end{center}
}

% ── Closing motivational rule + cross-promo footer (sheets only) ────────────────
% Usage: \SpeedClosing{"Quote text."}
\newcommand{\SpeedClosing}[1]{%
  \vspace{20pt}
  \begin{center}
  \rule{0.6\textwidth}{0.4pt}\\[4pt]
  {\small\itshape #1}\\[4pt]
  \rule{0.6\textwidth}{0.4pt}
  \end{center}
  \begin{center}
  {\small Found a mistake or want to contribute a sheet?}\\
  {\small Visit the open-source project at \href{https://github.com/The-CerealDev/speed-maths}{github.com/The-CerealDev/speed-maths}}
  \end{center}
}

% Editing THIS file changes the look of every sheet/answer across every pillar.
% ══════════════════════════════════════════════════════════════════════════════

\usepackage[a4paper, top=25mm, bottom=25mm, left=22mm, right=22mm]{geometry}
\usepackage{amsmath, amssymb}
\usepackage{enumitem}
\usepackage{booktabs}
\usepackage{array}
\usepackage{parskip}
\usepackage{microtype}
\usepackage{fancyhdr}
\usepackage{titlesec}
\usepackage{mdframed}
\usepackage{xcolor}
\usepackage[hidelinks]{hyperref}
\usepackage{graphicx}
\usepackage{tikz}
\usepackage{pgfplots}
\pgfplotsset{compat=1.18}

% ── Colours ───────────────────────────────────────────────────────────────────
\definecolor{darkgray}{gray}{0.15}
\definecolor{medgray}{gray}{0.45}
\definecolor{lightgray}{gray}{0.93}
\definecolor{boxgray}{gray}{0.88}

% ── Section formatting ──────────────────────────────────────────────────────────
\titleformat{\section}[block]
  {\normalfont\large\bfseries}{}
  {0pt}{}[\vspace{2pt}\hrule\vspace{6pt}]
\titlespacing*{\section}{0pt}{14pt}{4pt}

% ── List formatting ─────────────────────────────────────────────────────────────
\setlist[enumerate,1]{label=\textbf{\arabic*.}, leftmargin=2em, itemsep=10pt}

% ── Branding constants (edit here, not per-file) ────────────────────────────────
\newcommand{\SpeedExamLine}{\textbf{Speed Maths:} Competition \& Exam Prep}
\newcommand{\SpeedCredit}{Series by \href{https://www.linkedin.com/in/cerealdev/}{David Akinyele-Aje}}

% ── Header / footer ──────────────────────────────────────────────────────────────
% Usage (once, after % main (standalone preamble for editing)
% vim: nowrap: spell:
% ══════════════════════════════════════════════════════════════════════════════
% Speed Maths --- shared preamble
% Included by every sheet and answer file via:  \input{../../shared/preamble}
% Editing THIS file changes the look of every sheet/answer across every pillar.
% ══════════════════════════════════════════════════════════════════════════════

\usepackage[a4paper, top=25mm, bottom=25mm, left=22mm, right=22mm]{geometry}
\usepackage{amsmath, amssymb}
\usepackage{enumitem}
\usepackage{booktabs}
\usepackage{array}
\usepackage{parskip}
\usepackage{microtype}
\usepackage{fancyhdr}
\usepackage{titlesec}
\usepackage{mdframed}
\usepackage{xcolor}
\usepackage[hidelinks]{hyperref}
\usepackage{graphicx}
\usepackage{tikz}
\usepackage{pgfplots}
\pgfplotsset{compat=1.18}

% ── Colours ───────────────────────────────────────────────────────────────────
\definecolor{darkgray}{gray}{0.15}
\definecolor{medgray}{gray}{0.45}
\definecolor{lightgray}{gray}{0.93}
\definecolor{boxgray}{gray}{0.88}

% ── Section formatting ──────────────────────────────────────────────────────────
\titleformat{\section}[block]
  {\normalfont\large\bfseries}{}
  {0pt}{}[\vspace{2pt}\hrule\vspace{6pt}]
\titlespacing*{\section}{0pt}{14pt}{4pt}

% ── List formatting ─────────────────────────────────────────────────────────────
\setlist[enumerate,1]{label=\textbf{\arabic*.}, leftmargin=2em, itemsep=10pt}

% ── Branding constants (edit here, not per-file) ────────────────────────────────
\newcommand{\SpeedExamLine}{\textbf{Speed Maths:} Competition \& Exam Prep}
\newcommand{\SpeedCredit}{Series by \href{https://www.linkedin.com/in/cerealdev/}{David Akinyele-Aje}}

% ── Header / footer ──────────────────────────────────────────────────────────────
% Usage (once, after \input{../../shared/preamble} and before \begin{document}):
%   \SpeedHeader{Algebra}{3}
% First arg  = pillar name shown as "Speed <Pillar> --- Daily Drill"
% Second arg = worksheet number
\pagestyle{fancy}
\newcommand{\SpeedHeader}[2]{%
  \fancyhf{}%
  \renewcommand{\headrulewidth}{0.4pt}%
  \setlength{\headheight}{14pt}%
  \fancyhead[L]{\small\textbf{Speed #1 --- Daily Drill}}%
  \fancyhead[R]{\small Worksheet \##2}%
  \fancyfoot[L]{\small \SpeedExamLine}%
  \fancyfoot[C]{\small\thepage}%
  \fancyfoot[R]{\small \SpeedCredit}%
  \renewcommand{\footrulewidth}{0.4pt}%
}

% ── Title block (used at the top of every sheet AND answer file) ───────────────
% Usage: \SpeedTitleBlock{Daily Algebra Drill \#3}{Jane Doe}
% %% AGENTS: When generating a new sheet, ask the human contributor for the name they want credited in argument 2.
% For answer files, pass the "--- Answers \& Investigations" suffix in the arg.
\newcommand{\SpeedTitleBlock}[2]{%
  \begin{center}
    {\LARGE\bfseries #1}\\[4pt]
    {\normalsize\itshape Sheet by #2}\\[6pt]
    \hrule\vspace{4pt}
    \begin{tabular}{llcc}
      \toprule
      \textbf{Section} & \textbf{Topic} & \textbf{Questions} & \textbf{Target time}\\
      \midrule
      A & Rapid Recognition          & 10 & 2 min 30 sec \\
      B & Manipulation Drills        & 10 & 8 min        \\
      C & Substitution \& Structure  & 8  & 10 min       \\
      D & Challenge                  & 5  & 15 min       \\
      \bottomrule
    \end{tabular}
    \vspace{4pt}
    \hrule
  \end{center}
}

% ── Sheet metadata: topic + the day's new tools ────────────────────────────────
% Usage (immediately after \SpeedTitleBlock, sheets only):
%   \SpeedMeta{Expanding, factorising, and the identity toolkit}
%             {difference of squares, sum/product form, completing the square}
%
% Arg 1 = the sheet's topic in one line. The website lists this instead of
%         "Sheet 03", so write the thing a reader would actually search for.
% Arg 2 = comma-separated, at most five: the tools this sheet introduces that
%         earlier sheets in the pillar did not. These become the website's tags.
%         Leave out anything the reader already met — the tags are meant to say
%         what is new today, not to summarise the sheet.
%
% Both are parsed straight out of this file by tools/build_website.py, so the
% sheet stays the single source of truth and the site cannot drift from it.
%
% This renders NOTHING. It is a declaration for the build script, not a piece of
% the sheet's layout: adding it to a sheet must not change that sheet's PDF, so
% the 25 existing sheets can be annotated without a single one of them being
% reissued. If the toolkit should also appear on the page, that is a separate
% decision about the sheet design — make it deliberately, here, not by leaving
% this macro printing something nobody asked it to.
\newcommand{\SpeedMeta}[2]{}

% ── Answer / Method / Investigate-further boxes (answer files only) ────────────
\newmdenv[
  backgroundcolor=lightgray,
  linecolor=medgray,
  linewidth=0.5pt,
  leftmargin=0pt, rightmargin=0pt,
  innerleftmargin=8pt, innerrightmargin=8pt,
  innertopmargin=5pt, innerbottommargin=5pt,
  skipabove=4pt, skipbelow=4pt
]{ansbox}

\newmdenv[
  backgroundcolor=white,
  linecolor=darkgray,
  linewidth=0.8pt,
  leftline=true, rightline=false, topline=false, bottomline=false,
  leftmargin=0pt, rightmargin=0pt,
  innerleftmargin=8pt, innerrightmargin=6pt,
  innertopmargin=4pt, innerbottommargin=4pt,
  skipabove=4pt, skipbelow=6pt
]{invbox}

\newcommand{\ans}[1]{%
  \begin{ansbox}
    \textbf{Answer:} #1
  \end{ansbox}}

\newcommand{\method}[1]{%
  {\small\textbf{Method:} #1}\par}

\newcommand{\inv}[1]{%
  \begin{invbox}
    \textbf{\small Investigate further:} {\small #1}
  \end{invbox}}

% ── Misc helpers ────────────────────────────────────────────────────────────────
\newcommand{\qn}[1]{\textbf{#1.}\;}

% Mistake-linking: attach a Top 5 Patterns / Common Traps bullet to the question
% label(s) in this sheet where it actually shows up, e.g. \seealso{B6, D2}
\newcommand{\seealso}[1]{\hfill\textit{\footnotesize(see #1)}}

% ── Graph options macro ─────────────────────────────────────────────────────────
% Usage: \GraphOptions{tikz code for A}{tikz code for B}{tikz code for C}{tikz code for D}
\newcommand{\GraphOptions}[4]{%
  \begin{center}
    \begin{tabular}{cc}
      \textbf{A)} & \textbf{B)} \\
      #1 & #2 \\[10pt]
      \textbf{C)} & \textbf{D)} \\
      #3 & #4
    \end{tabular}
  \end{center}
}

% ── Closing motivational rule + cross-promo footer (sheets only) ────────────────
% Usage: \SpeedClosing{"Quote text."}
\newcommand{\SpeedClosing}[1]{%
  \vspace{20pt}
  \begin{center}
  \rule{0.6\textwidth}{0.4pt}\\[4pt]
  {\small\itshape #1}\\[4pt]
  \rule{0.6\textwidth}{0.4pt}
  \end{center}
  \begin{center}
  {\small Found a mistake or want to contribute a sheet?}\\
  {\small Visit the open-source project at \href{https://github.com/The-CerealDev/speed-maths}{github.com/The-CerealDev/speed-maths}}
  \end{center}
}
 and before \begin{document}):
%   \SpeedHeader{Algebra}{3}
% First arg  = pillar name shown as "Speed <Pillar> --- Daily Drill"
% Second arg = worksheet number
\pagestyle{fancy}
\newcommand{\SpeedHeader}[2]{%
  \fancyhf{}%
  \renewcommand{\headrulewidth}{0.4pt}%
  \setlength{\headheight}{14pt}%
  \fancyhead[L]{\small\textbf{Speed #1 --- Daily Drill}}%
  \fancyhead[R]{\small Worksheet \##2}%
  \fancyfoot[L]{\small \SpeedExamLine}%
  \fancyfoot[C]{\small\thepage}%
  \fancyfoot[R]{\small \SpeedCredit}%
  \renewcommand{\footrulewidth}{0.4pt}%
}

% ── Title block (used at the top of every sheet AND answer file) ───────────────
% Usage: \SpeedTitleBlock{Daily Algebra Drill \#3}{Jane Doe}
% %% AGENTS: When generating a new sheet, ask the human contributor for the name they want credited in argument 2.
% For answer files, pass the "--- Answers \& Investigations" suffix in the arg.
\newcommand{\SpeedTitleBlock}[2]{%
  \begin{center}
    {\LARGE\bfseries #1}\\[4pt]
    {\normalsize\itshape Sheet by #2}\\[6pt]
    \hrule\vspace{4pt}
    \begin{tabular}{llcc}
      \toprule
      \textbf{Section} & \textbf{Topic} & \textbf{Questions} & \textbf{Target time}\\
      \midrule
      A & Rapid Recognition          & 10 & 2 min 30 sec \\
      B & Manipulation Drills        & 10 & 8 min        \\
      C & Substitution \& Structure  & 8  & 10 min       \\
      D & Challenge                  & 5  & 15 min       \\
      \bottomrule
    \end{tabular}
    \vspace{4pt}
    \hrule
  \end{center}
}

% ── Sheet metadata: topic + the day's new tools ────────────────────────────────
% Usage (immediately after \SpeedTitleBlock, sheets only):
%   \SpeedMeta{Expanding, factorising, and the identity toolkit}
%             {difference of squares, sum/product form, completing the square}
%
% Arg 1 = the sheet's topic in one line. The website lists this instead of
%         "Sheet 03", so write the thing a reader would actually search for.
% Arg 2 = comma-separated, at most five: the tools this sheet introduces that
%         earlier sheets in the pillar did not. These become the website's tags.
%         Leave out anything the reader already met — the tags are meant to say
%         what is new today, not to summarise the sheet.
%
% Both are parsed straight out of this file by tools/build_website.py, so the
% sheet stays the single source of truth and the site cannot drift from it.
%
% This renders NOTHING. It is a declaration for the build script, not a piece of
% the sheet's layout: adding it to a sheet must not change that sheet's PDF, so
% the 25 existing sheets can be annotated without a single one of them being
% reissued. If the toolkit should also appear on the page, that is a separate
% decision about the sheet design — make it deliberately, here, not by leaving
% this macro printing something nobody asked it to.
\newcommand{\SpeedMeta}[2]{}

% ── Answer / Method / Investigate-further boxes (answer files only) ────────────
\newmdenv[
  backgroundcolor=lightgray,
  linecolor=medgray,
  linewidth=0.5pt,
  leftmargin=0pt, rightmargin=0pt,
  innerleftmargin=8pt, innerrightmargin=8pt,
  innertopmargin=5pt, innerbottommargin=5pt,
  skipabove=4pt, skipbelow=4pt
]{ansbox}

\newmdenv[
  backgroundcolor=white,
  linecolor=darkgray,
  linewidth=0.8pt,
  leftline=true, rightline=false, topline=false, bottomline=false,
  leftmargin=0pt, rightmargin=0pt,
  innerleftmargin=8pt, innerrightmargin=6pt,
  innertopmargin=4pt, innerbottommargin=4pt,
  skipabove=4pt, skipbelow=6pt
]{invbox}

\newcommand{\ans}[1]{%
  \begin{ansbox}
    \textbf{Answer:} #1
  \end{ansbox}}

\newcommand{\method}[1]{%
  {\small\textbf{Method:} #1}\par}

\newcommand{\inv}[1]{%
  \begin{invbox}
    \textbf{\small Investigate further:} {\small #1}
  \end{invbox}}

% ── Misc helpers ────────────────────────────────────────────────────────────────
\newcommand{\qn}[1]{\textbf{#1.}\;}

% Mistake-linking: attach a Top 5 Patterns / Common Traps bullet to the question
% label(s) in this sheet where it actually shows up, e.g. \seealso{B6, D2}
\newcommand{\seealso}[1]{\hfill\textit{\footnotesize(see #1)}}

% ── Graph options macro ─────────────────────────────────────────────────────────
% Usage: \GraphOptions{tikz code for A}{tikz code for B}{tikz code for C}{tikz code for D}
\newcommand{\GraphOptions}[4]{%
  \begin{center}
    \begin{tabular}{cc}
      \textbf{A)} & \textbf{B)} \\
      #1 & #2 \\[10pt]
      \textbf{C)} & \textbf{D)} \\
      #3 & #4
    \end{tabular}
  \end{center}
}

% ── Closing motivational rule + cross-promo footer (sheets only) ────────────────
% Usage: \SpeedClosing{"Quote text."}
\newcommand{\SpeedClosing}[1]{%
  \vspace{20pt}
  \begin{center}
  \rule{0.6\textwidth}{0.4pt}\\[4pt]
  {\small\itshape #1}\\[4pt]
  \rule{0.6\textwidth}{0.4pt}
  \end{center}
  \begin{center}
  {\small Found a mistake or want to contribute a sheet?}\\
  {\small Visit the open-source project at \href{https://github.com/The-CerealDev/speed-maths}{github.com/The-CerealDev/speed-maths}}
  \end{center}
}
 and before \begin{document}):
%   \SpeedHeader{Algebra}{3}
% First arg  = pillar name shown as "Speed <Pillar> --- Daily Drill"
% Second arg = worksheet number
\pagestyle{fancy}
\newcommand{\SpeedHeader}[2]{%
  \fancyhf{}%
  \renewcommand{\headrulewidth}{0.4pt}%
  \setlength{\headheight}{14pt}%
  \fancyhead[L]{\small\textbf{Speed #1 --- Daily Drill}}%
  \fancyhead[R]{\small Worksheet \##2}%
  \fancyfoot[L]{\small \SpeedExamLine}%
  \fancyfoot[C]{\small\thepage}%
  \fancyfoot[R]{\small \SpeedCredit}%
  \renewcommand{\footrulewidth}{0.4pt}%
}

% ── Title block (used at the top of every sheet AND answer file) ───────────────
% Usage: \SpeedTitleBlock{Daily Algebra Drill \#3}{Jane Doe}
% %% AGENTS: When generating a new sheet, ask the human contributor for the name they want credited in argument 2.
% For answer files, pass the "--- Answers \& Investigations" suffix in the arg.
\newcommand{\SpeedTitleBlock}[2]{%
  \begin{center}
    {\LARGE\bfseries #1}\\[4pt]
    {\normalsize\itshape Sheet by #2}\\[6pt]
    \hrule\vspace{4pt}
    \begin{tabular}{llcc}
      \toprule
      \textbf{Section} & \textbf{Topic} & \textbf{Questions} & \textbf{Target time}\\
      \midrule
      A & Rapid Recognition          & 10 & 2 min 30 sec \\
      B & Manipulation Drills        & 10 & 8 min        \\
      C & Substitution \& Structure  & 8  & 10 min       \\
      D & Challenge                  & 5  & 15 min       \\
      \bottomrule
    \end{tabular}
    \vspace{4pt}
    \hrule
  \end{center}
}

% ── Sheet metadata: topic + the day's new tools ────────────────────────────────
% Usage (immediately after \SpeedTitleBlock, sheets only):
%   \SpeedMeta{Expanding, factorising, and the identity toolkit}
%             {difference of squares, sum/product form, completing the square}
%
% Arg 1 = the sheet's topic in one line. The website lists this instead of
%         "Sheet 03", so write the thing a reader would actually search for.
% Arg 2 = comma-separated, at most five: the tools this sheet introduces that
%         earlier sheets in the pillar did not. These become the website's tags.
%         Leave out anything the reader already met — the tags are meant to say
%         what is new today, not to summarise the sheet.
%
% Both are parsed straight out of this file by tools/build_website.py, so the
% sheet stays the single source of truth and the site cannot drift from it.
%
% This renders NOTHING. It is a declaration for the build script, not a piece of
% the sheet's layout: adding it to a sheet must not change that sheet's PDF, so
% the 25 existing sheets can be annotated without a single one of them being
% reissued. If the toolkit should also appear on the page, that is a separate
% decision about the sheet design — make it deliberately, here, not by leaving
% this macro printing something nobody asked it to.
\newcommand{\SpeedMeta}[2]{}

% ── Answer / Method / Investigate-further boxes (answer files only) ────────────
\newmdenv[
  backgroundcolor=lightgray,
  linecolor=medgray,
  linewidth=0.5pt,
  leftmargin=0pt, rightmargin=0pt,
  innerleftmargin=8pt, innerrightmargin=8pt,
  innertopmargin=5pt, innerbottommargin=5pt,
  skipabove=4pt, skipbelow=4pt
]{ansbox}

\newmdenv[
  backgroundcolor=white,
  linecolor=darkgray,
  linewidth=0.8pt,
  leftline=true, rightline=false, topline=false, bottomline=false,
  leftmargin=0pt, rightmargin=0pt,
  innerleftmargin=8pt, innerrightmargin=6pt,
  innertopmargin=4pt, innerbottommargin=4pt,
  skipabove=4pt, skipbelow=6pt
]{invbox}

\newcommand{\ans}[1]{%
  \begin{ansbox}
    \textbf{Answer:} #1
  \end{ansbox}}

\newcommand{\method}[1]{%
  {\small\textbf{Method:} #1}\par}

\newcommand{\inv}[1]{%
  \begin{invbox}
    \textbf{\small Investigate further:} {\small #1}
  \end{invbox}}

% ── Misc helpers ────────────────────────────────────────────────────────────────
\newcommand{\qn}[1]{\textbf{#1.}\;}

% Mistake-linking: attach a Top 5 Patterns / Common Traps bullet to the question
% label(s) in this sheet where it actually shows up, e.g. \seealso{B6, D2}
\newcommand{\seealso}[1]{\hfill\textit{\footnotesize(see #1)}}

% ── Graph options macro ─────────────────────────────────────────────────────────
% Usage: \GraphOptions{tikz code for A}{tikz code for B}{tikz code for C}{tikz code for D}
\newcommand{\GraphOptions}[4]{%
  \begin{center}
    \begin{tabular}{cc}
      \textbf{A)} & \textbf{B)} \\
      #1 & #2 \\[10pt]
      \textbf{C)} & \textbf{D)} \\
      #3 & #4
    \end{tabular}
  \end{center}
}

% ── Closing motivational rule + cross-promo footer (sheets only) ────────────────
% Usage: \SpeedClosing{"Quote text."}
\newcommand{\SpeedClosing}[1]{%
  \vspace{20pt}
  \begin{center}
  \rule{0.6\textwidth}{0.4pt}\\[4pt]
  {\small\itshape #1}\\[4pt]
  \rule{0.6\textwidth}{0.4pt}
  \end{center}
  \begin{center}
  {\small Found a mistake or want to contribute a sheet?}\\
  {\small Visit the open-source project at \href{https://github.com/The-CerealDev/speed-maths}{github.com/The-CerealDev/speed-maths}}
  \end{center}
}

\SpeedHeader{Logic}{5}

\begin{document}

\SpeedTitleBlock{Daily Logic Drill \#5 --- Answers \& Investigations}{David Akinyele-Aje}
\noindent\textit{New toolkit today: Spot the Flaw $\cdot$ Common Fallacies $\cdot$ Proof Critique $\cdot$ Valid vs.\ Invalid Arguments.}


\section*{Section A}

\begin{enumerate}

%── A1 ──────────────────────────────────────────────────────────────────────────
\item A student solves $\sqrt{x-2} = -3$ by squaring both sides to obtain $x-2 = 9 \implies x=11$. Does $x=11$ satisfy the original equation?

\ans{False.}
\method{Substituting $x=11$ gives $\sqrt{11-2} = \sqrt{9} = 3 \neq -3$. Squaring both sides introduces an extraneous solution where the two sides originally had opposite signs. Thus $x=11$ does not satisfy the equation.}
\inv{Under what condition on the right-hand side does $\sqrt{f(x)} = c$ have any real solutions?}

%── A2 ──────────────────────────────────────────────────────────────────────────
\item A student simplifies $\dfrac{x^2-4}{x-2} = 4 \implies x+2 = 4 \implies x=2$. Is $x=2$ a valid solution to the original equation?

\ans{False.}
\method{At $x=2$, the denominator $x-2$ equals $0$, so the expression $\frac{x^2-4}{x-2}$ is undefined. The cancellation $\frac{x^2-4}{x-2} = x+2$ is valid only for $x \neq 2$. Thus $x=2$ is not in the domain of the equation, leaving it with no solutions.}
\inv{What type of discontinuity does the function $g(x) = \frac{x^2-4}{x-2}$ have at $x=2$?}

%── A3 ──────────────────────────────────────────────────────────────────────────
\item In an attempted proof that $1=2$, a student starts with $a=b$, obtains $(a-b)(a+b) = b(a-b)$, and concludes $a+b = b$. Which algebraic operation is invalid?

\ans{Dividing by $a-b$ (division by zero is invalid).}
\method{Since $a=b$, the factor $a-b$ is equal to $0$. Dividing both sides of $(a-b)(a+b) = b(a-b)$ by $a-b$ constitutes division by zero, which is undefined.}
\inv{Why does dividing by zero allow one to prove any arbitrary equality such as $1=2$?}

%── A4 ──────────────────────────────────────────────────────────────────────────
\item A student claims that for all real $x$, $x^2 > 4 \implies x > 2$. Give a negative integer counterexample to this deduction.

\ans{$-3$}
\method{For $x = -3$, $x^2 = 9 > 4$ holds, but $-3 > 2$ is false. Any $x \le -2$ disproves the claim.}
\inv{What is the correct equivalence for $x^2 > 4$?}

%── A5 ──────────────────────────────────────────────────────────────────────────
\item A student solves $\log(x) + \log(x-3) = 1$ and finds candidate solutions $x = -2$ and $x = 5$. Which candidate is extraneous?

\ans{$-2$}
\method{For $x = -2$, neither $\log(-2)$ nor $\log(-5)$ is defined for real logarithms, which require strictly positive arguments. Only $x = 5$ is valid: $\log(5) + \log(2) = \log(10) = 1$.}
\inv{Why does combining $\log(a) + \log(b) = \log(ab)$ expand the domain of definition?}

%── A6 ──────────────────────────────────────────────────────────────────────────
\item A student conjectures: ``If $f'(c) = 0$, then $f$ has a local maximum or minimum at $x=c$.'' Give a polynomial function $f(x)$ with $c=0$ that disproves this claim.

\ans{$f(x) = x^3$ (an inflection point, no local extremum).}
\method{For $f(x) = x^3$, $f'(x) = 3x^2 \implies f'(0) = 0$. However, $x=0$ is a stationary point of inflection: $f(x) < 0$ for $x < 0$ and $f(x) > 0$ for $x > 0$, so $f$ has neither a local maximum nor a local minimum at $x=0$.}
\inv{What condition on higher derivatives determines whether a stationary point is a local extremum?}

%── A7 ──────────────────────────────────────────────────────────────────────────
\item A student asserts: ``$|a| < |b| \implies a < b$.'' Provide a counterexample pair of integers $(a, b)$.

\ans{$(1, -2)$}
\method{For $a = 1$ and $b = -2$, we have $|a| = 1 < 2 = |b|$, but $1 < -2$ is false.}
\inv{When does $|a| < |b| \implies a < b$ hold?}

%── A8 ──────────────────────────────────────────────────────────────────────────
\item A student concludes that because $n^2$ is divisible by $4$, $n$ must be divisible by $4$. Find the smallest positive integer counterexample $n$.

\ans{$2$}
\method{For $n = 2$, $n^2 = 4$, which is divisible by $4$, but $2$ is not divisible by $4$. The smallest counterexample is $n=2$.}
\inv{If $p$ is prime and $p \mid n^2$, must $p \mid n$?}

%── A9 ──────────────────────────────────────────────────────────────────────────
\item A student claims: ``For all real numbers $a, b$, if $a < b$ then $a^2 < b^2$.'' Give a pair of integers $(a, b)$ with $a < b$ and $a^2 > b^2$.

\ans{(-5, 2)}
\method{For $a = -5$ and $b = 2$, $-5 < 2$ holds, but $a^2 = 25 > 4 = b^2$.}
\inv{Under what condition on $a$ and $b$ does $a < b \implies a^2 < b^2$ hold?}

%── A10 ─────────────────────────────────────────────────────────────────────────
\item A student asserts: ``If $A \subseteq B \cup C$, then $A \subseteq B$ or $A \subseteq C$.'' Give a set $A$ of two elements and sets $B=\{1\}, C=\{2\}$ that disprove this claim.

\ans{$A = \{1, 2\}$}
\method{Let $A = \{1, 2\}$. Then $A \subseteq \{1\} \cup \{2\} = \{1, 2\}$ holds, but $A \not\subseteq \{1\}$ and $A \not\subseteq \{2\}$.}
\inv{Is it true that if $x \in B \cup C$, then $x \in B$ or $x \in C$?}

\end{enumerate}

\section*{Section B}

\begin{enumerate}

%── B1 ──────────────────────────────────────────────────────────────────────────
\item A student attempts to prove: ``For all integers $n$, if $n$ is odd, then $n^2$ is odd.''\\
Line 1: Suppose $n$ is odd, so $n=2k+1$ for some integer $k$.\\
Line 2: Then $n^2=4k^2+4k+1=2(2k^2+2k)+1$.\\
Line 3: Since $2k^2+2k$ is an integer, $n^2$ is odd.\\
Is this proof valid? Justify your answer.

\ans{Yes, fully valid.}
\method{Every line follows correctly: the algebra in Line 2 is correct ($(2k+1)^2=4k^2+4k+1$), and Line 3 correctly identifies the result as $2\times(\text{integer})+1$, the definition of odd. There is no flaw.}
\inv{Is this a direct proof or a proof by contrapositive of some related statement?}

%── B2 ──────────────────────────────────────────────────────────────────────────
\item Find the flaw: ``Claim: If $x^2=3x$, then $x=3$.'' Proof: Divide both sides of $x^2=3x$ by $x$ to get $x=3$.

\ans{Dividing by $x$ is invalid when $x=0$.}
\method{$x^2=3x \iff x(x-3)=0 \iff x=0$ or $x=3$. Dividing by $x$ silently assumes $x\neq0$ and loses the solution $x=0$, which also satisfies the original equation.}
\inv{Rewrite the claim correctly, including both solutions.}

%── B3 ──────────────────────────────────────────────────────────────────────────
\item Find the flaw: ``Claim: The only real solution to $\sqrt{x+7}=x-5$ is $x=9$.'' Proof: Squaring both sides gives $x+7=x^2-10x+25$, so $x^2-11x+18=0$, so $(x-9)(x-2)=0$, giving $x=9$ or $x=2$.

\ans{$x=2$ is an extraneous root introduced by squaring; it must be checked against and rejected from the original equation.}
\method{Checking $x=2$ in the original equation: $\sqrt{2+7}=\sqrt9=3$, but $x-5=2-5=-3$. Since $3\neq-3$, $x=2$ does not satisfy the original equation. Squaring introduced this extraneous solution; only $x=9$ (which does check: $\sqrt{16}=4=9-5$) is valid.}
\inv{Why does squaring specifically introduce solutions where the two sides had opposite sign?}

%── B4 ──────────────────────────────────────────────────────────────────────────
\item A student argues: For every integer $n$, $6\left(\dfrac{4n+1}{2}-\dfrac{n-3}{2}\right)$ is a multiple of $3$. \textit{\small(after TMUA 2020 Paper 2 Q3)}\\
Line 1: $6\left(\dfrac{4n+1}{2}-\dfrac{n-3}{2}\right)=3\left(2\left(\dfrac{4n+1}{2}-\dfrac{n-3}{2}\right)\right)$.\\
Line 2: $=3(4n+1-n-3)$.\\
Line 3: $=3(3n-2)$.\\
Line 4: which is always a multiple of $3$.\\
On which line does the argument first contain an error, and what exactly is that error?

\ans{Line 2}
\method{Correctly distributing, $2\left(\frac{4n+1}{2}-\frac{n-3}{2}\right)=(4n+1)-(n-3)=4n+1-n+3=3n+4$, not $4n+1-n-3$ as written. The sign of $-3$ was not flipped when distributing the subtraction across the bracket $(n-3)$. (Interestingly, since Line 1's outer factor of $3$ is untouched by this error, the final conclusion --- a multiple of $3$ --- remains true regardless; the argument is still invalid even though its conclusion happens to be correct.)}
\inv{Compute the genuinely correct simplified form, $3(3n+4)$, and confirm it is still a multiple of $3$.}

%── B5 ──────────────────────────────────────────────────────────────────────────
\item Find the flaw: ``Claim: For every integer $n$, $n^2+n+2$ is even.'' Proof: If $n$ is even, $n=2k$, so $n^2+n+2=4k^2+2k+2=2(2k^2+k+1)$, which is even. Hence $n^2+n+2$ is even for every integer $n$.

\ans{The proof only checks the even case; the odd case is never addressed.}
\method{To prove a claim for ``every integer $n$'', both parities must be checked. (The claim does happen to also hold for odd $n$: $n=2k+1 \implies n^2+n+2=4k^2+4k+1+2k+1+2=4k^2+6k+4=2(2k^2+3k+2)$, even --- but the proof as given never establishes this.)}
\inv{Complete the missing odd case explicitly.}

%── B6 ──────────────────────────────────────────────────────────────────────────
\item Find the flaw: ``Claim: For real $x$, $x^2=9 \iff x=3$.'' Proof: If $x=3$, then $x^2=9$. This proves the biconditional.

\ans{Only the $\Longleftarrow$ direction ($x=3\implies x^2=9$) was shown; the $\Longrightarrow$ direction ($x^2=9\implies x=3$) is false, since $x=-3$ also satisfies $x^2=9$.}
\method{A biconditional requires both directions to be proved. The correct statement is $x^2=9 \iff x=3 \text{ or } x=-3$.}
\inv{What is the correct biconditional relating $x^2=9$ to a single equation in $x$?}

%── B7 ──────────────────────────────────────────────────────────────────────────
\item Find the flaw: ``Claim: For every positive integer $n$, the $n$-th Fibonacci number $F_n$ (with $F_1=F_2=1$) is either equal to $1$ or prime.'' Proof: $F_3=2$, $F_4=3$, $F_5=5$ --- all prime. Since the pattern holds for $n=3,4,5$, the claim is proved. What is the smallest counterexample?

\ans{Checking finitely many cases does not prove a universal claim; $F_6=8=2^3$ is the smallest counterexample.}
\method{$F_1=1,F_2=1,F_3=2,F_4=3,F_5=5,F_6=8$. $F_6=8$ is neither equal to $1$ nor prime ($8=2\times4$), disproving the claim. Three verified cases give no logical guarantee about the fourth.}
\inv{Is $F_n$ prime for infinitely many $n$? This remains an open problem in general.}

%── B8 ──────────────────────────────────────────────────────────────────────────
\item A student wants to prove: ``If $n$ is not divisible by $3$, then $n^2$ is not divisible by $3$.'' Instead, they show: ``If $n$ is divisible by $3$, then $n^2$ is divisible by $3$.'' Explain why this does not prove the original claim.

\ans{The statement shown is neither the original claim nor its contrapositive --- it proves nothing about the original.}
\method{Original claim: $\neg A\implies\neg B$ where $A$: ``$3\mid n$'', $B$: ``$3\mid n^2$''. Its valid contrapositive is $B\implies A$ (``if $3\mid n^2$ then $3\mid n$''). The student instead proved $A\implies B$, a different (though also true) statement, which establishes neither the original claim nor its contrapositive.}
\inv{State the actual contrapositive of the original claim, and outline how you would prove it.}

%── B9 ──────────────────────────────────────────────────────────────────────────
\item Find the flaw: ``Claim: For all integers $a,b$, if $a$ and $b$ are both even, then $a+b$ is even.'' Proof: Since $a+b$ is even, $a+b=2m$ for some integer $m$. Since $a,b$ are even, $a=2j,b=2k$, so $2j+2k=2m$, i.e.\ $j+k=m$, which is true since $m$ is an integer. Hence $a+b$ is even.

\ans{Circular reasoning: the proof begins by assuming $a+b$ is even (the very thing to be proved) instead of deriving it.}
\method{A correct proof starts from $a=2j,b=2k$ and derives $a+b=2j+2k=2(j+k)$, which is even because $j+k$ is an integer --- it never needs to assume $a+b=2m$ up front.}
\inv{Rewrite this as a correct, non-circular direct proof.}

%── B10 ─────────────────────────────────────────────────────────────────────────
\item A student attempts to solve $\log_2(x-1) + \log_2(x+1) = 3$:\\
Line 1: $\log_2((x-1)(x+1)) = 3$.\\
Line 2: $(x-1)(x+1) = 2^3 = 8$.\\
Line 3: $x^2 - 1 = 8 \implies x^2 = 9$.\\
Line 4: $x = 3$ or $x = -3$.\\
Line 5: Hence the solution set is $\{-3, 3\}$.\\
On which line does the first error occur, and what is the true solution set? \textit{\small(after TMUA 2017 Paper 2 Q13)}

\ans{Line 5 (only $x=3$ is valid in the domain).}
\method{Lines 1--4 correctly deduce $x = \pm 3$. However, in the original equation $\log_2(x-1) + \log_2(x+1) = 3$, the terms require $x-1 > 0$ and $x+1 > 0$, so $x > 1$. The candidate $x = -3$ makes the arguments negative ($\log_2(-4)$ and $\log_2(-2)$), which are undefined. Line 5 fails to verify the domain restriction; the true solution set is $\{3\}$.}
\inv{Why does combining $\log(a) + \log(b) = \log(ab)$ expand the domain to include points where both $a$ and $b$ are negative?}

\end{enumerate}

\section*{Section C}

\begin{enumerate}

%── C1 ──────────────────────────────────────────────────────────────────────────
\item Consider the claim: ``For every positive integer $n$, $2n^2+2n+1$ is prime.'' \textit{\small(after TMUA 2022 Paper 2 Q7)}

\ans{D}
\method{Lines I--III are correct: the discriminant of $2x^2+2x+1$ is indeed $-4$, so the polynomial has no real (let alone integer linear) factorisation. But Line IV commits a category error: a polynomial being irreducible over $\mathbb{Z}[x]$ says nothing about whether its integer OUTPUT VALUES are prime. At $n=3$: $2(9)+6+1=25=5^2$, composite, disproving the original claim.}
\inv{Find another value of $n$ for which $2n^2+2n+1$ is composite.}

%── C2 ──────────────────────────────────────────────────────────────────────────
\item Let $P$: ``$x^2-5x+6=0$'' and $Q$: ``$x=2$''. A student argues $Q\implies P$ is true and concludes $P\iff Q$.

\ans{B}
\method{$x^2-5x+6=(x-2)(x-3)=0$, so $P$ is satisfied by $x=2$ or $x=3$. $Q\implies P$ is indeed true, but $P\implies Q$ is false since $x=3$ satisfies $P$ without satisfying $Q$. Asserting $P\implies Q$ merely because the statements ``seem related'' is the flaw.}
\inv{What is the correct relationship (necessary/sufficient) between $P$ and $Q$?}

%── C3 ──────────────────────────────────────────────────────────────────────────
\item A student attempts to prove: ``For every integer $n$, $n^2+n$ is divisible by $2$'', checking only the even case.

\ans{B}
\method{$n^2+n=n(n+1)$ is a product of two consecutive integers, so it is even regardless of the parity of $n$ --- but the given proof only demonstrates this for even $n$ and never addresses odd $n$, so as written it is incomplete.}
\inv{Complete the missing odd case: for $n=2k+1$, show $n^2+n$ is even.}

%── C4 ──────────────────────────────────────────────────────────────────────────
\item A student attempts to prove: ``No three consecutive odd positive integers can all be prime.'' \textit{\small(after TMUA 2023 Paper 2 Q4)}

\ans{C}
\method{Lines I and II are correct: among $n-2,n,n+2$, exactly one is always divisible by $3$. Line III is the flaw --- being a positive multiple of $3$ only implies compositeness if the multiple is \emph{greater than} $3$; the multiple could equal $3$ itself, which is prime. Exactly this happens at $n=5$: the three consecutive odd integers $3,5,7$ are all prime.}
\inv{Are there any other counterexamples to the original claim besides $3,5,7$?}

%── C5 ──────────────────────────────────────────────────────────────────────────
\item A student ``proves'' $\forall n\,\exists m\,(m>n)$ and concludes $\exists M\,\forall n\,(M>n)$.

\ans{B}
\method{Swapping the order of a universal and existential quantifier is not a valid logical step in general. $\forall n\,\exists m\,(m>n)$ says every $n$ has SOME larger $m$ (true, e.g.\ $m=n+1$); $\exists M\,\forall n\,(M>n)$ says a SINGLE $M$ exceeds every $n$ simultaneously, which is false (no integer is larger than every integer).}
\inv{Give a general example of a true $\forall\exists$ statement whose $\exists\forall$ swap is false.}

%── C6 ──────────────────────────────────────────────────────────────────────────
\item A student ``confirms'' the converse of ``rhombus $\implies$ perpendicular diagonals'' using only the example of a square.

\ans{B}
\method{One example (the square) supports the converse but cannot prove a universal statement. A kite has perpendicular diagonals but is generally not a rhombus, disproving the converse.}
\inv{Sketch a kite that is not a rhombus, and confirm its diagonals are perpendicular.}

%── C7 ──────────────────────────────────────────────────────────────────────────
\item A student disproves $(*)$: ``For every pair of real numbers $x,y$, $|x+y|<|x|+|y|$'' using the pair $x=-2,y=-5$. \textit{\small(after TMUA 2020 Paper 2 Q8)}

\ans{D}
\method{This is a trap: the student's reasoning is actually completely correct. $(*)$ as stated uses strict $<$, and at $x=-2,y=-5$ (both negative, hence same sign), $|x+y|=7=|x|+|y|$, so $7<7$ is false --- a fully valid single counterexample to a universal claim. No further pairs need be checked to reject $(*)$.}
\inv{For which pairs $(x,y)$ does $|x+y|=|x|+|y|$ hold with equality?}

%── C8 ──────────────────────────────────────────────────────────────────────────
\item A student ``verifies'' ``four equal sides $\iff$ square'' by checking only that every square has four equal sides.

\ans{B}
\method{This checks only necessity (square $\implies$ four equal sides). Sufficiency (four equal sides $\implies$ square) is false: a non-square rhombus has four equal sides but is not a square.}
\inv{What additional condition, combined with ``four equal sides'', would give a genuine necessary and sufficient characterisation of a square?}

\end{enumerate}

\section*{Section D}

\begin{enumerate}

%── D1 ──────────────────────────────────────────────────────────────────────────
\item A flawed attempt to prove: if $s^2\geq4p$ then there exist real $u,v$ with $u+v=s,uv=p$. \textit{\small(after TMUA 2021 Paper 2 Q11)}

\ans{B}
\method{Every algebraic step (II, III) is correct, but the logical direction is backward. The student assumed $u,v$ exist and derived $s^2\geq4p$ from that assumption --- this proves the CONVERSE of the theorem, not the theorem itself. A correct proof constructs $u,v$ directly: let $u,v$ be the two roots of $t^2-st+p=0$; since the discriminant $s^2-4p\geq0$, these roots are real, and by Vieta's formulas $u+v=s,uv=p$.}
\inv{Why does the discriminant condition $s^2-4p\geq0$ guarantee the roots of $t^2-st+p=0$ are real?}

%── D2 ──────────────────────────────────────────────────────────────────────────
\item A student claims ``$n^2-n+11$ prime $\implies n^2+n+11$ prime'' for all positive $n$, verified for $n=1,\dots,9$.

\ans{(a) Finite verification is not a universal proof. (b) $n=10$.}
\method{(a) Checking $9$ cases gives no guarantee about $n=10$ or beyond --- a single unchecked value could be a counterexample.\\
(b) At $n=10$: $n^2-n+11=100-10+11=101$, which is prime (premise true). But $n^2+n+11=100+10+11=121=11^2$, which is NOT prime (conclusion false). So the implication fails at $n=10$, the smallest such value since $n=1,\dots,9$ were verified to hold.}
\inv{Verify that $n^2-n+11$ is prime for $n=1,\dots,9$ by computing the first few values.}

%── D3 ──────────────────────────────────────────────────────────────────────────
\item Evaluate a contradiction-style proof that ``$ab$ even $\implies a$ even or $b$ even''.

\ans{D) I and II only}
\method{The algebra $(2j+1)(2k+1)=4jk+2j+2k+1=2(2jk+j+k)+1$ is correct, so III is false. Statement I is true: assuming the negation ($a,b$ both odd) and deriving a contradiction ($ab$ odd, contradicting $ab$ even) is a valid contradiction proof. Statement II is also true: the derivation ``$a,b$ both odd $\implies ab$ odd'' is exactly the contrapositive of the original claim (contrapositive of $ab\text{ even}\implies(a\text{ even or }b\text{ even})$ is $\neg(a\text{ even or }b\text{ even})\implies\neg(ab\text{ even})$, i.e.\ ``$a,b$ both odd $\implies ab$ odd''), proved here directly without needing the ``suppose for contradiction'' framing at all.}
\inv{Rewrite this proof in pure contrapositive form, removing the contradiction framing entirely.}

%── D4 ──────────────────────────────────────────────────────────────────────────
\item Evaluate the ``every odd prime $>3$ has the form $6k\pm1$'' argument.

\ans{B}
\method{Every line is correct: $6k,6k+2,6k+4$ are even (Line 2, correctly dismissed); $6k+3=3(2k+1)$ is a multiple of $3$ exceeding $3$ whenever $k\geq1$, and the case $k=0$ giving exactly $3$ is excluded by the claim's own hypothesis ``greater than $3$'' (Line 3, correctly handled). The proof only ever claimed to show that primes $>3$ MUST have the form $6k\pm1$ (necessity) --- it never claimed every number of that form IS prime (sufficiency), and the original claim, read carefully, only asserts the former. So there is no actual gap.}
\inv{Give an example of a number of the form $6k+1$ that is not prime, to confirm sufficiency genuinely fails.}

%── D5 ──────────────────────────────────────────────────────────────────────────
\item Evaluate: ``For all $n\geq2$, $n^2-1$ is composite'', proved via $n^2-1=(n-1)(n+1)$.

\ans{(a) Line 3. (b) $n=2$; and $n=2$ is the unique exception.}
\method{(a) Line 3 is the gap: a product of two integers each ``at least $1$'' is not necessarily composite --- compositeness requires BOTH factors to be at least $2$. Line 2 only guarantees $n-1\geq1$, not $n-1\geq2$.\\
(b) At $n=2$: $n^2-1=3=1\times3$. The factor of $1$ does not make $3$ composite --- $3$ is prime, so the conclusion genuinely fails here. \emph{Proof that $n=2$ is the only failure}: for every $n\geq3$, $n-1\geq2$ and $n+1\geq4$, so $n^2-1=(n-1)(n+1)$ is a product of two integers both at least $2$, which is precisely the definition of composite --- no case analysis or further checking is needed once $n\geq3$, since the factorisation itself, together with $n-1\geq2$, is a fully general and complete argument covering every such $n$ simultaneously.}
\inv{Why does the same style of argument fail to show $n^2-4$ is composite for all $n\geq3$? (Check $n=3$: $n^2-4=5$.)}

\end{enumerate}

%═══════════════════════════════════════════════════════════════════════════════
\section*{Top 5 Patterns Today}
%═══════════════════════════════════════════════════════════════════════════════

\begin{enumerate}[leftmargin=2em, itemsep=4pt]
  \item \textbf{The Converse Error (Affirming the Consequent).}
        Arguing that $P\Rightarrow Q$ and observing $Q$ implies $P$ is the classic converse fallacy; observing a necessary condition holding never proves the sufficient premise. \seealso{B6, C2, C6, C8, D1}
  \item \textbf{The Inverse Error (Denying the Antecedent).}
        Arguing that $P\Rightarrow Q$ and concluding $\neg P\Rightarrow\neg Q$ confuses an implication with its inverse; multiple distinct conditions often produce the exact same consequence. \seealso{B8, D3, C2}
  \item \textbf{Extraneous Roots from Non-Reversible Operations.}
        Squaring both sides of $\sqrt{f(x)}=g(x)$ or applying logarithm product rules introduces false solutions where the original expressions were undefined or opposite in sign. \seealso{A1, A2, A5, B3, B10}
  \item \textbf{Division by Zero and Root Loss.}
        Dividing both sides of an equation by a variable expression (e.g.\ $x^2=3x\Rightarrow x=3$ or dividing by $a-b$) implicitly assumes that factor is nonzero, erasing valid roots or fabricating false contradictions. \seealso{A2, A3, B2, D5}
  \item \textbf{Incomplete Case Analysis and Boundary Gaps.}
        Proving a claim only for even $n$, or ignoring boundary values where factors equal $1$ ($n=2$ in $(n-1)(n+1)$), leaves unexamined cases that often harbour the unique counterexample. \seealso{B5, B7, C3, D2, D5}
\end{enumerate}

%═══════════════════════════════════════════════════════════════════════════════
\section*{Common Traps to Avoid}
%═══════════════════════════════════════════════════════════════════════════════

\begin{enumerate}[leftmargin=2em, itemsep=4pt]
  \item \textbf{The Fallacy Fallacy (Refuting Argument vs Refuting Claim).}
        Demonstrating that a proof contains an invalid step shows only that the proof fails, not that the conclusion itself is false; a true proposition can easily be given an invalid proof. \seealso{B4, C1, D2}
  \item \textbf{Losing Valid Roots via Direct Variable Division.}
        Cancelling a factor $x$ from $x^2=kx$ discards the root $x=0$; always factorise into $x(x-k)=0$ rather than dividing through by unknown quantities. \seealso{A2, A3, B2, D5}
  \item \textbf{Accepting Extraneous Solutions After Squaring or Logs.}
        Squaring an equation maps both positive and negative values to positive numbers; all solutions found after squaring or logarithm products must be verified against original domains. \seealso{A1, A5, B3, B10}
  \item \textbf{Assuming ``Checking Cases'' Equals Proof for Universals.}
        Verifying that an algebraic property holds for several small integers ($n=1,2,3$) is merely inductive evidence; universal claims fail if even a single case or parity class is omitted. \seealso{B5, B7, C3, D2}
  \item \textbf{Circular Reasoning (Begging the Question).}
        Assuming the conclusion in order to simplify an intermediate expression produces an argument that appears valid on the surface but is logically vacuous. \seealso{B9, C4}
\end{enumerate}

\SpeedClosing{``Speed comes from recognition, not from rushing. Every shortcut is a pattern you have internalised.''}

\end{document}
